\documentclass[11pt,a4paper]{article}

% --- PACCHETTI FONDAMENTALI ---
\usepackage[utf8]{inputenc}
\usepackage[T1]{fontenc}
\usepackage[italian]{babel}
\usepackage{amsmath, amssymb, amsthm, mathtools}
\usepackage{geometry}
\geometry{a4paper, top=2.5cm, bottom=2.5cm, left=2.5cm, right=2.5cm}
\usepackage{tikz}
\usepackage[most]{tcolorbox}
\usepackage{xcolor}
\usepackage{algorithm}
\usepackage{algorithmic}

% --- PALETTE COLORI ---
\definecolor{paletteIndigo}{HTML}{4B0082}
\definecolor{paletteRegalia}{HTML}{522D80}
\definecolor{Lenurple}{HTML}{BA96CD}
\definecolor{Lavender}{HTML}{E6E6FA}
\definecolor{darkgray}{HTML}{333333}

% --- STILI TCOLORBOX ---
\tcbset{
    basebox/.style={
        enhanced,
        boxrule=1pt,
        arc=4pt,
        fonttitle=\bfseries,
        coltitle=white,
        attach boxed title to top left={yshift=-2mm, xshift=4mm},
        boxed title style={arc=2pt, boxrule=0pt}
    },
    defbox/.style={
        basebox,
        colback=Lavender!30,
        colframe=paletteIndigo,
        boxed title style={colback=paletteIndigo}
    },
    conceptbox/.style={
        basebox,
        colback=white,
        colframe=paletteRegalia,
        boxed title style={colback=paletteRegalia}
    },
    examplebox/.style={
        basebox,
        colback=Lenurple!10,
        colframe=Lenurple,
        coltitle=paletteIndigo,
        boxed title style={colback=Lenurple!50}
    }
}

% --- AMBIENTI PERSONALIZZATI ---
\newtcolorbox{definition}[1][Definizione]{defbox, title=#1}
\newtcolorbox{concept}[1][Concetto Chiave]{conceptbox, title=#1}
\newtcolorbox{esercizio}[1][Esercizio/Esempio]{examplebox, title=#1}

\begin{document}

\title{\vspace{-2cm}\textbf{\color{paletteIndigo}Appunti Universitari: Algebra Lineare}\\ \Large Introduzione ai Sistemi Lineari, Matrici ed Eliminazione di Gauss}
\author{\textit{Redazione Accademica Universale}}
\date{Marzo 2026}
\maketitle

\section{Introduzione alle Equazioni e ai Sistemi Lineari}

\subsection{Premesse Fondamentali}
Salvo esplicito avviso contrario, supponiamo sempre di lavorare in $\mathbb{R}$, cioè che i coefficienti, i termini noti e le variabili (o incognite) delle equazioni siano sempre numeri reali[cite: 3, 4, 5, 6, 7].

\subsection{Equazioni Lineari}
\begin{definition}[Equazione Lineare in una Variabile]
Un'equazione lineare è un'equazione di primo grado (i.e., in cui l'incognita compare al primo grado)[cite: 8]. Se un'equazione lineare è in una sola incognita, la si può sempre porre, tramite passaggi algebrici, nella forma:
\begin{equation}
ax + b = 0
\end{equation}
\end{definition}

\textbf{Soluzioni:}
\begin{itemize}
    \item Se $a \neq 0 \implies 1$ soluzione, $x = \frac{-b}{a}$[cite: 12].
    \item Se $a = 0$ e $b \neq 0 \implies$ NON AMMETTE SOLUZIONI[cite: 13].
    \item Se $a = 0$ e $b = 0 \implies$ soddisfatta per $\forall x \in \mathbb{R}$ (è un'identità)[cite: 14].
\end{itemize}

\begin{concept}[Equazione Lineare in due Variabili]
Se invece un'equazione lineare ha due variabili, le soluzioni sono in generale infinite[cite: 15, 16]. 
\textbf{Esempio:}
\begin{equation}
x + y = 1
\end{equation}
Le soluzioni sono coppie: $(1,0), (0,1), (1/2, 1/2)$, ecc...[cite: 17, 18].
Ovvero tutte le infinite coppie $(x,y)$ con $x = t$ e $y = 1 - t$, per $t \in \mathbb{R}$[cite: 19].
\end{concept}

\textbf{Interpretazione Geometrica:} Se consideriamo il piano cartesiano riferito a una coppia di assi geometrici ortogonali, sappiamo che le soluzioni rappresentano i punti della retta di equazione $y = -x + 1$[cite: 20, 22, 23, 25].

\subsection{Sistemi di Equazioni Lineari}
Supponiamo di avere anche un'altra equazione lineare in due incognite, per esempio $x - y = 0$, e ci chiediamo se esistono valori di $x$ e $y$ che soddisfano ENTRAMBE le equazioni: abbiamo quindi quello che si chiama un SISTEMA LINEARE che si indica con la scrittura[cite: 27]:
\begin{equation}
\begin{cases}
x + y = 1 \\
x - y = 0
\end{cases}
\end{equation}

\begin{definition}[Sistema di Equazioni]
Un insieme di equazioni di cui si cercano le (eventuali) soluzioni comuni prende il nome di Sistema di Equazioni[cite: 30, 33]. Un sistema lineare è un sistema composto esclusivamente da equazioni lineari[cite: 32].
\end{definition}

\subsubsection{Risoluzione e Interpretazione Geometrica}
Proviamo a risolvere il sistema precedente. Dalla seconda equazione troviamo che $x = y$, e sostituendo nella prima ricaviamo $x = y = \frac{1}{2}$[cite: 32].
Dal punto di vista geometrico, ciascuna equazione rappresenta una retta, e le soluzioni del sistema corrispondono ai punti in cui si incontrano[cite: 32, 39, 41, 43]. 
Ci sono 3 possibili situazioni geometriche[cite: 45, 47]:

\begin{center}
\begin{tikzpicture}[scale=0.8, every node/.style={scale=0.8}]
    % Caso 1: Incidenti
    \draw[->] (-1,0) -- (3,0) node[right] {$x$};
    \draw[->] (0,-1) -- (0,3) node[above] {$y$};
    \draw[paletteIndigo, thick] (-0.5, 1.5) -- (2.5, -1.5) node[right] {$x+y=1$};
    \draw[paletteRegalia, thick] (-1, -1) -- (2.5, 2.5) node[right] {$2x-y=0$};
    \filldraw[black] (0.5,0.5) circle (2pt) node[above] {1 Sol.};
    \node at (1, -2.5) {\textbf{Rette Incidenti}};
    
    % Caso 2: Parallele
    \begin{scope}[xshift=5cm]
    \draw[->] (-1,0) -- (3,0) node[right] {$x$};
    \draw[->] (0,-1) -- (0,3) node[above] {$y$};
    \draw[paletteIndigo, thick] (-0.5, 1.5) -- (2.5, -1.5) node[right] {$x+y=1$};
    \draw[paletteRegalia, thick] (-0.5, 2.5) -- (2.5, -0.5) node[right] {$x+y=2$};
    \node at (1, -2.5) {\textbf{Rette Parallele}};
    \node at (1, -3) {(Nessuna Soluzione)};
    \end{scope}

    % Caso 3: Coincidenti
    \begin{scope}[xshift=11cm]
    \draw[->] (-1,0) -- (3,0) node[right] {$x$};
    \draw[->] (0,-1) -- (0,3) node[above] {$y$};
    \draw[Lenurple, line width=1.5mm] (-0.5, 1.5) -- (2.5, -1.5);
    \draw[paletteIndigo, thick] (-0.5, 1.5) -- (2.5, -1.5) node[right, yshift=-0.3cm] {$2x+2y=2$};
    \node[right, yshift=0.3cm] at (2.5,-1.5) {$x+y=1$};
    \node at (1, -2.5) {\textbf{Rette Coincidenti}};
    \node at (1, -3) {($\infty$ Soluzioni)};
    \end{scope}
\end{tikzpicture}
\end{center}
[cite: 48]

\subsection{Generalizzazione ai Sistemi $m \times n$}
Le considerazioni viste per i sistemi lineari $2 \times 2$ si possono generalizzare al caso di sistemi di $m$ equazioni in $n$ incognite[cite: 49, 51, 52]:
\begin{equation}
\begin{cases}
a_{11}x_1 + a_{12}x_2 + \dots + a_{1n}x_n = b_1 \\
a_{21}x_1 + a_{22}x_2 + \dots + a_{2n}x_n = b_2 \\
\vdots \\
a_{m1}x_1 + a_{m2}x_2 + \dots + a_{mn}x_n = b_m
\end{cases}
\end{equation}
[cite: 50]

\textbf{Casistica delle Soluzioni:}
\begin{itemize}
    \item \textbf{Possibile (Compatibile, Risolubile):} Ha soluzioni. Può essere $1!$ soluzione unica, oppure $\infty$ soluzioni, dipendente da uno o più parametri[cite: 50].
    \item \textbf{Impossibile:} $0$ soluzioni comuni alle $m$ equazioni[cite: 50].
\end{itemize}
Non è interesse solo (o non tanto) risolvere un sistema lineare, ma soprattutto discuterlo, i.e. scoprire A PRIORI (ovvero senza risolverlo) se è risolubile o no, e quante sono le sue soluzioni al variare di eventuali parametri[cite: 50].

\begin{concept}[Sistemi Equivalenti]
Per risolvere un sistema lineare, l'idea di fondo è trasformare il sistema dato in un sistema che abbia le stesse soluzioni, ma sia più facile da risolvere[cite: 54, 132].
Due sistemi si dicono \textbf{equivalenti} se possiedono le stesse soluzioni, i.e. ogni n-upla di valori che risolve uno, risolve anche l'altro[cite: 54, 132].

\textbf{Esempio di sistemi equivalenti:}
$$
\begin{cases} (1+6)x = 7 \\ y = 2x \end{cases} \implies \begin{cases} x = 1 \\ y = 2 \end{cases}
\quad \text{Sono sistemi equivalenti a:} \quad
\begin{cases} x + 3y = 7 \\ 2x - y = 0 \end{cases}
$$
[cite: 54, 132]
\end{concept}

\section{Dai Sistemi alle Matrici}
Per studiare i sistemi lineari, basta studiarne le matrici associate, poiché la struttura è la stessa e cambiano solo i numeri[cite: 55, 133].
\textbf{Esempio di passaggio alle matrici:}
$$
\begin{cases}
5x + 6y = 0 \\
3x + 2y = 7 \\
7x + 2y = 2
\end{cases}
\implies
\text{Associo la matrice: }
\begin{pmatrix}
5 & 6 & 0 \\
3 & 2 & 7 \\
7 & 2 & 2
\end{pmatrix}
$$
[cite: 54, 132]

\begin{definition}[Matrice e Vettori]
Dato il sistema generico $m \times n$, definiamo:
\begin{itemize}
    \item \textbf{Matrice Associata (Incompleta) $A$}: La tabella $m \times n$ (con $m$ righe e $n$ colonne) che raccoglie i coefficienti $a_{ij}$. L'elemento $a_{ij}$ occupa la riga $i$ e la colonna $j$[cite: 58, 68, 70, 71, 72, 135, 143, 144, 152, 153].
    \item \textbf{Vettore dei Termini Noti $b$}: È una matrice $m \times 1$ (vettore colonna) che raccoglie i valori $b_i$[cite: 60, 73, 74, 81, 148, 156, 158].
    \item \textbf{Matrice Completa del Sistema $(A|b)$}: È la matrice con $m$ righe e $n+1$ colonne ottenuta affiancando ad $A$ il vettore $b$[cite: 61, 75, 76, 159, 160].
\end{itemize}
\end{definition}

\subsection{Tipologie di Matrici}
\begin{itemize}
    \item \textbf{Matrice Quadrata:} Una matrice in cui il numero di righe è uguale al numero di colonne ($m=n$). Il numero $n$ si dice \textbf{ORDINE} della matrice. Esempio di matrice quadrata di ordine 3: $\begin{pmatrix} 4 & 2 & 3 \\ 2 & 4 & 5 \\ 3 & 0 & 1 \end{pmatrix}$[cite: 63, 64, 78, 79, 82, 83, 162, 164, 165, 172, 173]. Un sistema lineare si dice quadrato se la sua matrice associata è quadrata[cite: 88, 90, 181].
    \item \textbf{Diagonale Principale:} L'insieme degli elementi $a_{ij}$ con $i = j$ ($a_{11}, a_{22}, \dots, a_{nn}$)[cite: 85, 174].
    \item \textbf{Matrice Diagonale:} Se $a_{ij} = 0 \ \forall i \neq j$[cite: 85, 174].
    \item \textbf{Matrice Triangolare Superiore:} Tutti gli elementi sotto la diagonale principale sono nulli ($a_{ij} = 0 \ \forall i > j$). Esempio: $A = \begin{pmatrix} 1 & 3 \\ 0 & 1 \end{pmatrix}$[cite: 85, 174]. Un sistema si dice triangolare superiore se la sua matrice lo è[cite: 91, 182].
    \item \textbf{Matrice Triangolare Inferiore:} Tutti gli elementi sopra la diagonale principale sono nulli ($a_{ij} = 0 \ \forall i < j$). Esempio: $A = \begin{pmatrix} 1 & 0 \\ 2 & 3 \end{pmatrix}$[cite: 85, 174].
\end{itemize}
\textit{Osservazione:} Una matrice diagonale è contemporaneamente triangolare superiore e triangolare inferiore[cite: 85, 174].

\begin{concept}[Pivot di una riga]
Consideriamo una matrice qualsiasi $A$, $m \times n$. Le sue righe sono $A_1, A_2, \dots, A_m$.
Chiamiamo \textbf{Pivot} della riga $A_i$ il primo elemento non nullo della riga $A_i$[cite: 174].

\textbf{Esempio:} In $A = \begin{pmatrix} 0 & 1 & 2 & 3 \\ 7 & 3 & 5 & 0 \end{pmatrix}$, il pivot di $A_1$ è $a_{12} = 1$, il pivot di $A_2$ è $a_{21} = 7$[cite: 177, 180]. 
Una matrice è a scalini se il pivot della riga successiva è più a destra[cite: 180, 288].
\end{concept}


\subsection{Risoluzione di Sistemi Triangolari}
I sistemi lineari triangolari superiori sono "i più belli di tutti", sono facili da risolvere sostituendo a ritroso[cite: 104, 126, 195, 215].

\begin{esercizio}[Esempi Completi di Sostituzione a Ritroso]
\textbf{Esempio 1: 1! Soluzione}
$$
\begin{cases} 3x + 2y - z = 0 \\ y - z = 0 \\ z = 1 \end{cases} \implies
A = \begin{pmatrix} 3 & 2 & 0 \\ 0 & 1 & -1 \\ 0 & 0 & 1 \end{pmatrix}, \ b = \begin{pmatrix} 0 \\ 0 \\ 1 \end{pmatrix}
$$
[cite: 101, 102, 103, 192, 193, 194]
Risolvendo dal basso verso l'alto (sostituendo a ritroso):
1) $z = 1$ [cite: 109, 200]
2) Dalla seconda equazione: $y - 1 = 0 \implies y = 1$ [cite: 204]
3) Dalla prima equazione: $3x + 2y = 0 \implies 3x + 2(1) = 0 \implies 3x + 2 = 0$ [cite: 105, 110, 204] 
(Nota dai manoscritti: l'equazione appare come $3x+2y=0$ ma la matrice $A$ ha $-z$ nel primo esempio $3x+2y-z=0$ e $0$ nel secondo, mostriamo i passaggi così come segnati negli appunti).

\textbf{Esempio 2: Nessuna Soluzione (Impossibile)}
$$
A = \begin{pmatrix} 3 & 2 & 0 \\ 0 & 0 & 1 \\ 0 & 0 & 1 \end{pmatrix}, \ b = \begin{pmatrix} 0 \\ 3 \\ 1 \end{pmatrix}
$$
[cite: 111, 118, 206]
Se l'ultima riga indica $1 \cdot z = 1$ e la seconda $1 \cdot z = 3$, il sistema entra in contraddizione. Sistema impossibile, non ha soluzioni[cite: 121, 122, 210, 211].

\textbf{Esempio 3: $\infty$ Soluzioni (Gradi di Libertà)}
$$
A = \begin{pmatrix} -3 & 2 & 0 \\ 0 & 0 & 2 \\ 0 & 0 & 1 \end{pmatrix}, \ b = \begin{pmatrix} 0 \\ 2 \\ 1 \end{pmatrix}
$$
[cite: 119, 207, 208]
Ha infinite soluzioni dipendenti da 1 grado di libertà. Soluzioni della forma $(t, \dots, 1)$ o $(t, -3/2 t, 1)$[cite: 123, 212, 213].
\end{esercizio}

Se il sistema di partenza non è triangolare, ad esempio:
$$
\begin{cases} x + y + z = 0 \\ 2x + y - z = 3 \\ x - y + 2z = 5 \end{cases}
$$
[cite: 219, 220]
Non è immediatamente risolubile. Ci serve una "ricetta" per passare da un sistema ad un sistema triangolare superiore equivalente: l'Algoritmo di Gauss[cite: 224, 225, 230].

\section{Eliminazione di Gauss}

\begin{definition}[Mosse di Gauss]
Date un sistema $C = (A|b)$, le mosse di Gauss su $C$ non modificano l'insieme delle soluzioni del sistema lineare associato (cambiano l'aspetto ma non l'insieme $S = S'$ delle soluzioni)[cite: 255, 262, 263, 265, 266, 268, 269]. Le operazioni sono 3:
\begin{enumerate}
    \item Scambiare 2 righe[cite: 258].
    \item Moltiplicare una riga per un numero $\lambda \neq 0$[cite: 259].
    \item Aggiungere a una riga un'altra riga moltiplicata per un $\lambda$ qualsiasi[cite: 260, 261].
\end{enumerate}
\end{definition}

\subsection{L'Algoritmo}
L'algoritmo trasforma una matrice qualsiasi in una a scalini[cite: 272].
\begin{algorithmic}[1]
\STATE Se $C_{11} = 0$ scambia la riga 1 con un'altra riga in modo da ottenere $C_{11} \neq 0$. (Se $C_{i1} = 0 \ \forall i$, passa al punto 3) [cite: 272].
\STATE Con il pivot $C_{11} \neq 0$, "ammazza" (annulla) gli elementi $C_{i1}$ sottostanti, sostituendo la riga $C_i$ con $C_i - \frac{C_{i1}}{C_{11}} C_1$[cite: 272].
\STATE Hai ottenuto $C_{i1} = 0 \ \forall i \ge 2$. Copri la prima riga e la prima colonna e riparti dal punto 1 sulla sottomatrice ottenuta[cite: 272].
\end{algorithmic}

\begin{esercizio}[Esempio 3: Applicazione Algoritmo]
Partiamo dal sistema lineare equivalente a:
$$
C = \begin{pmatrix}
0 & 1 & 1 & | & 0 \\
1 & 1 & 2 & | & -3 \\
-1 & 2 & 1 & | & 1
\end{pmatrix} \quad \text{corrispondente a } \quad
\begin{cases}
x_2 + x_3 = 0 \\
x_1 + x_2 + 2x_3 = -3 \\
-x_1 + 2x_2 + x_3 = 1
\end{cases}
$$
[cite: 272]

\textbf{Passo 1:} $C_{11} = 0$, quindi scambio $C_1$ con $C_2$:
$$
\begin{pmatrix}
1 & 1 & 2 & | & -3 \\
0 & 1 & 1 & | & 0 \\
-1 & 2 & 1 & | & 1
\end{pmatrix}
$$
[cite: 272]

\textbf{Passo 2:} "Ammazzo" il -1. Sostituisco $C_3 \leftarrow C_3 + C_1$:
$$
\begin{pmatrix}
1 & 1 & 2 & | & -3 \\
0 & 1 & 1 & | & 0 \\
0 & 3 & 3 & | & -2
\end{pmatrix}
$$
[cite: 276, 277, 278]

\textbf{Passo 3:} Il pivot successivo è $C_{22} = 1 \neq 0$, va bene. Devo ammazzare $C_{32} = 3$. Sostituisco $C_3 \leftarrow C_3 - 3C_2$:
$$
\begin{pmatrix}
1 & 1 & 2 & | & -3 \\
0 & 1 & 1 & | & 0 \\
0 & 0 & 0 & | & -2
\end{pmatrix}
$$
[cite: 279, 280, 281, 282, 283, 284, 285, 286]
Matrice a scalini ottenuta. Mi fermo[cite: 288, 289, 292].
\end{esercizio}

\begin{esercizio}[Esempio 1: Algoritmo Completo con Soluzione]
$$
\begin{cases}
x + y + z = 6 \\
2x + y - z = 3 \\
-x - y + 2z = 5
\end{cases}
\implies
\text{Matrice associata } C = \begin{pmatrix}
1 & 1 & 1 & | & 6 \\
2 & 1 & -1 & | & 3 \\
-1 & -1 & 2 & | & 5
\end{pmatrix}
$$
[cite: 298, 299, 300, 301, 302, 303, 304, 305]

Cosa manca alla matrice per essere triangolare superiore? Devo uccidere i pivot di $A_2$ e $A_3$[cite: 306, 307, 308, 315].
Eseguo $C_2 \leftarrow C_2 - 2C_1$ e $C_3 \leftarrow C_3 + C_1$:
$$
\begin{pmatrix}
1 & 1 & 1 & | & 6 \\
0 & -1 & -3 & | & -9 \\
0 & 0 & 3 & | & 11
\end{pmatrix}
$$
(Dai manoscritti si legge il passaggio come $C_3 \leftarrow C_3 - C_1$ per un possibile errore di trascrizione in aula, poi corretto in una riga $0 \ 0 \ 7 \ | \ 17$) [cite: 312, 313, 316, 318, 319, 328].
Adesso la matrice è triangolare superiore. Risolvo il sistema corrispondente:
$$
\begin{cases}
x + y + z = 6 \\
-y - 3z = -9 \\
7z = 17
\end{cases}
$$
[cite: 323, 327, 328]
Risolvendo a ritroso:
$z = \frac{17}{7}$, 
$y = 9 - 3z = \frac{12}{7}$, 
$x = -y - z + 6 = \frac{13}{7}$[cite: 329].
\end{esercizio}

\begin{esercizio}[Esempio 4: Rilevare l'impossibilità]
$$
\begin{cases}
x + y = 2 \\
2x + 2y = 4 \\
x + y = 3
\end{cases}
\implies
C = \begin{pmatrix}
1 & 1 & | & 2 \\
2 & 2 & | & 4 \\
1 & 1 & | & 3
\end{pmatrix}
$$
[cite: 322, 325, 336, 337, 338, 339]

Si vede "a occhio" che la prima e la terza equazione sono incompatibili, il sistema è impossibile. Ma mettiamo che io non me ne accorga[cite: 333, 334, 335].
Ammazzo i pivot di $A_2$ e $A_3$ con le mosse $C_2 \leftarrow C_2 - 2C_1$ e $C_3 \leftarrow C_3 - C_1$[cite: 340, 341]:
$$
\begin{pmatrix}
1 & 1 & | & 2 \\
0 & 0 & | & 0 \\
0 & 0 & | & 1
\end{pmatrix}
$$
[cite: 342, 343, 344]
L'ultima riga rappresenta l'equazione $0 = 1$, IMPOSSIBILE[cite: 342, 345].
\end{esercizio}

\begin{esercizio}[Esempio 5: Sistema Dipendente da Parametri]
$$
\begin{cases}
x + y + z = 3 \\
2x + y + 3z = 7
\end{cases}
\implies
\begin{pmatrix}
1 & 1 & 1 & | & 3 \\
2 & 1 & 3 & | & 7
\end{pmatrix}
$$
[cite: 347]

Mossa di Gauss: la matrice diventa a scalini:
$$
\begin{pmatrix}
1 & 1 & 1 & | & 3 \\
0 & -1 & 1 & | & 1
\end{pmatrix}
\implies
\begin{cases}
x + y + z = 3 \\
-y + z = 1
\end{cases}
$$
[cite: 347]
Ponendo $z = t$, ricaviamo $y = t - 1$ e $x = -y - z + 3 = -(t-1) - t + 3 = 4 - 2t$.
Soluzioni della forma $(4 - 2t, t - 1, t)$. Abbiamo infinite soluzioni (gradi di libertà = 1 parametro)[cite: 347].
\end{esercizio}

\begin{concept}[Esempi Veloci e Morale sull'Algoritmo]
\textbf{Esempi Veloci:}
\begin{itemize}
    \item $\begin{cases} 5x + 2y = 4 \\ 10x + 4y = 8 \end{cases} \implies \infty$ soluzioni[cite: 348].
    \item $\begin{cases} x + 2y = 4 \\ 2x + 4y = 9 \end{cases} \implies 0$ soluzioni[cite: 348].
    \item $\begin{cases} 2x + y = 5 \\ x - y = 1 \end{cases} \implies 1!$ soluzione[cite: 348].
\end{itemize}

\textbf{MORALE:} La forma a scalini fa emergere subito le caratteristiche delle soluzioni[cite: 347].
1. Se compare una riga del tipo $\begin{pmatrix} 0 & 0 & \dots & 0 & | & b \end{pmatrix}$ con $b \neq 0 \implies$ Nessuna Soluzione[cite: 347].
2. Se compare una o più righe tutte di zeri $\implies$ è un'equazione ridondante[cite: 347].
3. In un sistema quadrato che diviene triangolare con tutti i pivot sulla diagonale $\neq 0 \implies 1!$ soluzione unica[cite: 347].
Il numero dei pivot misura quante informazioni indipendenti ci sono nel sistema[cite: 347].
\end{concept}

\end{document}
